\section{Questão 03: Pós-Treino com QLoRA}
\label{sec:q3}

\subsection{Objetivo e configuração}

A Questão 03 repete o pós-treino da Questão 02 substituindo o ajuste fino completo
pela adaptação de baixo posto~\cite{hu2022lora} sobre o modelo quantizado em 4 bits
(QLoRA)~\cite{dettmers2023qlora}. Para uma comparação justa, foram reutilizados
exatamente os mesmos pares, a mesma divisão treino/avaliação, a mesma máscara de
perda e o mesmo tamanho de lote efetivo da Questão 02. A Tabela~\ref{tab:q3-hiper}
resume a configuração do adaptador.

\begin{table}[ht]
  \centering
  \caption{Configuração do QLoRA (Questão 03).}
  \label{tab:q3-hiper}
  \footnotesize
  \begin{tabularx}{0.92\textwidth}{@{}l X@{}}
    \toprule
    \textbf{Item} & \textbf{Valor} \\
    \midrule
    Quantização da base & 4 bits NF4 com dupla quantização, cômputo em \texttt{fp16} \\
    Posto $r$ / fator $\alpha$ & 16 / 32, \textit{dropout} 0,05 \\
    Módulos adaptados & projeções de atenção (\texttt{q}, \texttt{k}, \texttt{v}, \texttt{o}) \\
    Parâmetros treináveis & 2.162.688 (0,44\% do modelo) \\
    Taxa de aprendizado & $2 \times 10^{-4}$ com decaimento cosseno \\
    Épocas / lote efetivo & 3 / 16, mesmos dados e divisão da Questão 02 \\
    Tempo de treino / pico de VRAM & 1.079~s ($\approx$ 18~min) / 1,77~GB \\
    Tamanho do adaptador salvo & 8,7~MB \\
    \bottomrule
  \end{tabularx}
\end{table}

O plano original previa também a variante LoRA sem quantização, mas ela não pôde ser
executada: a biblioteca PEFT instalada no ambiente Kaggle rejeitou a criação dos
adaptadores sobre o modelo em \texttt{fp32} por incompatibilidade com a versão
antiga do pacote \texttt{torchao} pré-instalado na imagem. Trata-se de uma limitação
de infraestrutura, registrada no log da execução; como o enunciado solicita LoRA
e/ou QLoRA, a questão está atendida pela variante quantizada, cujo caminho de
criação dos adaptadores não passa pelo componente defeituoso.

\subsection{Efeito da quantização no modelo base}

Avaliar o modelo base nas duas precisões, sobre o mesmo conjunto de avaliação,
isola o custo da quantização: a perplexidade sobe de 14,57 (\texttt{fp32}) para
16,28 (4 bits), um acréscimo de 11,7\%, e a acurácia de tokens cai 2,5 pontos
percentuais (Figura~\ref{fig:q3-quant}). Esse é o preço de partida que o QLoRA paga
para treinar com uma fração da memória.

\begin{figure}[ht]
  \centering
  \includegraphics[width=\textwidth]{q3_quantizacao.png}
  \caption{Efeito da quantização NF4 sobre o modelo base, sem nenhum treino, no
  mesmo conjunto de avaliação da Questão 02.}
  \label{fig:q3-quant}
\end{figure}

\subsection{Resultados quantitativos}

Partindo do modelo quantizado (perplexidade 16,28), o QLoRA reduziu a perplexidade
sobre os tokens de resposta para 6,40 (queda de 60,7\%) e elevou a acurácia de
tokens de 48,78\% para 64,37\%. A Figura~\ref{fig:q3-metricas} posiciona esses
resultados ao lado dos da Questão 02, com os respectivos pontos de partida.

\begin{figure}[ht]
  \centering
  \includegraphics[width=\textwidth]{q3_metricas.png}
  \caption{Métricas finais por técnica, com os dois baselines explicitados: o SFT
  completo parte da base em \texttt{fp32} e o QLoRA parte da base quantizada em 4
  bits. Métricas sobre os tokens de resposta do mesmo conjunto de avaliação.}
  \label{fig:q3-metricas}
\end{figure}

\subsection{Comparação com o ajuste fino completo}

O resultado central da questão está nas Figuras~\ref{fig:q3-params}
e~\ref{fig:q3-painel}: o QLoRA atingiu perplexidade final \emph{menor} que a do
ajuste fino completo (6,40 contra 7,49) atualizando 228 vezes menos parâmetros, com
pico de memória 5,5 vezes menor (1,77~GB contra 9,77~GB) e um artefato final de
8,7~MB contra cerca de 1~GB do modelo completo. Duas ressalvas contextualizam a
comparação: as taxas de aprendizado diferem (valor usual de cada técnica) e o QLoRA
parte de um baseline pior por causa da quantização, o que torna sua vantagem final
ainda mais notável.

A explicação é visível nas curvas de validação: enquanto o SFT completo entrou em
\textit{overfitting} na terceira época (Seção~\ref{sec:q2}), a perda de validação do
QLoRA caiu monotonicamente (0,929; 0,881; 0,878). Com apenas 1.728 exemplos de
treino, o adaptador de posto 16 age como uma regularização implícita: há menos graus
de liberdade para decorar o conjunto de treino.

\begin{figure}[ht]
  \centering
  \includegraphics[width=0.9\textwidth]{q3_lora_params.png}
  \caption{Parâmetros atualizados por técnica, em escala logarítmica. A razão entre
  o SFT completo e o QLoRA do mesmo modelo é de 228 vezes.}
  \label{fig:q3-params}
\end{figure}

\begin{figure}[ht]
  \centering
  \includegraphics[width=0.94\textwidth]{q3_painel_comparativo.png}
  \caption{Custo e qualidade das três variantes de pós-treino executadas sobre os
  mesmos dados. O QLoRA domina o SFT completo em perplexidade e memória; o modelo de
  1,5 bilhão quantizado obtém o melhor resultado geral.}
  \label{fig:q3-painel}
\end{figure}

A avaliação qualitativa, porém, inverte o ranking e completa o quadro: nas mesmas
dez perguntas da Questão 02, o QLoRA de 0,5 bilhão acertou apenas 2 (contra 6 do SFT
completo), falhando principalmente nos fatos fechados sobre o dataset. A leitura
conjunta é que atualizar todos os pesos memoriza conhecimento factual novo com mais
força, enquanto o adaptador aprende bem o formato e a distribuição das respostas,
mas retém menos fatos em um modelo desse porte. Afirmar superioridade do QLoRA
apenas pela perplexidade seria, portanto, uma conclusão incompleta.

\subsection{Experimento adicional: QLoRA no Qwen2.5-1.5B}

Atendendo ao item opcional do enunciado, a mesma configuração de QLoRA foi aplicada
ao \texttt{Qwen2.5-1.5B}. O modelo maior partiu de um baseline quantizado bem melhor
(perplexidade 10,28, inferior até à base de 0,5 bilhão em \texttt{fp32}) e terminou
com os melhores números do pós-treino: perplexidade 4,75, acurácia de tokens de
68,44\% e 5 acertos qualitativos em 10, treinando 0,28\% dos parâmetros em 2,75~GB
de VRAM. Escala e adaptação eficiente se mostraram complementares: o 1,5B com QLoRA
aproxima-se da qualidade factual do SFT completo gastando uma fração dos recursos.
